\documentclass[12pt]{standalone}
\usepackage{fontspec}
\setsansfont{TeX Gyre Heros}
\usepackage{tikz}
\usetikzlibrary{arrows.meta}
% =====================================================================
% 统一符号层（2026-10-20）
% 用法：在 fig_<号>.tex 的 preamble 里 % =====================================================================
% 统一符号层（2026-10-20）
% 用法：在 fig_<号>.tex 的 preamble 里 % =====================================================================
% 统一符号层（2026-10-20）
% 用法：在 fig_<号>.tex 的 preamble 里 \input{../_symbols/symbols.tex}
% 依赖：fontspec（已 \setsansfont{TeX Gyre Heros}）、graphicx 不需要
% 规则见同目录《符号验收清单.md》
% =====================================================================

% ---- 字体角色（全部走 sans = TeX Gyre Heros）----
\newcommand{\SIt}[1]{\textsf{\textit{#1}}}          % 斜体变量
\newcommand{\SBf}[1]{\textsf{\textbf{#1}}}          % 粗体
\newcommand{\SBfIt}[1]{\textsf{\textbf{\textit{#1}}}} % 粗斜体变量
\newcommand{\SNum}[1]{\textsf{\textup{#1}}}         % 直立数字

% ---- 上下标：字号/位置写死，避免数学样式与手调混用 ----
% 审计结论（2026-10-20）：源图下标/底字高比在 0.77–0.92em，0.72em 偏小。
% 默认 0.80em；仍需按源图逐图标定（44.4/44.5 约 0.85–0.92，39/41/46 约 0.77–0.80）。
\newcommand{\FSup}[2]{#1\raisebox{0.30em}{\fontsize{0.72em}{0.82em}\selectfont #2}}
\newcommand{\FSub}[2]{#1\raisebox{-0.30em}{\fontsize{0.80em}{0.92em}\selectfont #2}}

% ---- 分数：数字直立、方头粗杆；杆宽 = max(分子宽,分母宽) + 2pt ----
% 杆厚 4pt（配合 20–26pt 标签）；小字号图可局部调 0.17em
% 2026-10-20 修正：原 \shortstack 负 kern 不随字号缩放（23pt 下总高仅 30pt、杆贴字），
% 改为 offinterlineskip + em 定间距：上间隙 0.12em、下间隙 0.13em
% （37.2 源图 23pt 实测四标签上/下间隙均值 2.6/4.0pt；取近平衡且总高保持 42–43pt）。
\newcommand{\FFrac}[2]{%
  \begingroup
  \setbox0=\hbox{#1}\setbox2=\hbox{#2}%
  \dimen0=\wd0 \ifdim\wd2>\dimen0 \dimen0=\wd2 \fi
  \advance\dimen0 by 2pt
  \leavevmode\vbox{\offinterlineskip
    \hbox to \dimen0{\hss\box0\hss}%
    \kern0.12em
    \hbox{\rule{\dimen0}{4pt}}%
    \kern0.13em
    \hbox to \dimen0{\hss\box2\hss}%
  }%
  \endgroup}

% ---- 真减号（审计：源图减号 15–18px 宽、3–5px 厚；Heros 常规字重偏细）----
% 可选参数 = 宽度（em）：正文 0.58；上标内用 \FMinus[0.72]
\newcommand{\FMinus}[1][0.58]{\rule[0.30em]{#1em}{0.13em}}

% ---- 直撇（审计：44.2 用弯引号、44.3 用直撇，必须统一为直楔撇）----
\newcommand{\FPrime}{\textsf{\textquotesingle}}

% ---- 乘号（审计：46.2 的 × 比源小 25–30%，标签内需按源校准）----
% 用 \FTimes 时在 node 内与底字同字号，并在目检时量宽高。


% ---- 斜杠分数（审计 A：35.1 的 r/3 用 38pt 斜杠手搓）----
\newcommand{\FSlashFrac}[2]{#1\raisebox{-0.18em}{\fontsize{1.35em}{1.35em}\selectfont\textsf{/}}\kern-0.06em #2}

% ---- 交集号 ∩（Heros 无此字，固定回退字号/基线，避免逐图不同）----
\newcommand{\FIntersect}{\raisebox{0.02em}{\fontsize{0.95em}{1em}\selectfont\mfallback{∩}}}

% ---- 括号/开区间（图形层，不是字形）----
% 一律 \draw ... arc[start angle=, end angle=, radius=]：
%   · 端点记号背离它所界的区间；
%   · 左端/右端分别 10× 特写核对；
%   · 闭区间压在轴上时竖脊右移半个线宽；
%   · 同一张图同类记号同一种写法。

% 依赖：fontspec（已 \setsansfont{TeX Gyre Heros}）、graphicx 不需要
% 规则见同目录《符号验收清单.md》
% =====================================================================

% ---- 字体角色（全部走 sans = TeX Gyre Heros）----
\newcommand{\SIt}[1]{\textsf{\textit{#1}}}          % 斜体变量
\newcommand{\SBf}[1]{\textsf{\textbf{#1}}}          % 粗体
\newcommand{\SBfIt}[1]{\textsf{\textbf{\textit{#1}}}} % 粗斜体变量
\newcommand{\SNum}[1]{\textsf{\textup{#1}}}         % 直立数字

% ---- 上下标：字号/位置写死，避免数学样式与手调混用 ----
% 审计结论（2026-10-20）：源图下标/底字高比在 0.77–0.92em，0.72em 偏小。
% 默认 0.80em；仍需按源图逐图标定（44.4/44.5 约 0.85–0.92，39/41/46 约 0.77–0.80）。
\newcommand{\FSup}[2]{#1\raisebox{0.30em}{\fontsize{0.72em}{0.82em}\selectfont #2}}
\newcommand{\FSub}[2]{#1\raisebox{-0.30em}{\fontsize{0.80em}{0.92em}\selectfont #2}}

% ---- 分数：数字直立、方头粗杆；杆宽 = max(分子宽,分母宽) + 2pt ----
% 杆厚 4pt（配合 20–26pt 标签）；小字号图可局部调 0.17em
% 2026-10-20 修正：原 \shortstack 负 kern 不随字号缩放（23pt 下总高仅 30pt、杆贴字），
% 改为 offinterlineskip + em 定间距：上间隙 0.12em、下间隙 0.13em
% （37.2 源图 23pt 实测四标签上/下间隙均值 2.6/4.0pt；取近平衡且总高保持 42–43pt）。
\newcommand{\FFrac}[2]{%
  \begingroup
  \setbox0=\hbox{#1}\setbox2=\hbox{#2}%
  \dimen0=\wd0 \ifdim\wd2>\dimen0 \dimen0=\wd2 \fi
  \advance\dimen0 by 2pt
  \leavevmode\vbox{\offinterlineskip
    \hbox to \dimen0{\hss\box0\hss}%
    \kern0.12em
    \hbox{\rule{\dimen0}{4pt}}%
    \kern0.13em
    \hbox to \dimen0{\hss\box2\hss}%
  }%
  \endgroup}

% ---- 真减号（审计：源图减号 15–18px 宽、3–5px 厚；Heros 常规字重偏细）----
% 可选参数 = 宽度（em）：正文 0.58；上标内用 \FMinus[0.72]
\newcommand{\FMinus}[1][0.58]{\rule[0.30em]{#1em}{0.13em}}

% ---- 直撇（审计：44.2 用弯引号、44.3 用直撇，必须统一为直楔撇）----
\newcommand{\FPrime}{\textsf{\textquotesingle}}

% ---- 乘号（审计：46.2 的 × 比源小 25–30%，标签内需按源校准）----
% 用 \FTimes 时在 node 内与底字同字号，并在目检时量宽高。


% ---- 斜杠分数（审计 A：35.1 的 r/3 用 38pt 斜杠手搓）----
\newcommand{\FSlashFrac}[2]{#1\raisebox{-0.18em}{\fontsize{1.35em}{1.35em}\selectfont\textsf{/}}\kern-0.06em #2}

% ---- 交集号 ∩（Heros 无此字，固定回退字号/基线，避免逐图不同）----
\newcommand{\FIntersect}{\raisebox{0.02em}{\fontsize{0.95em}{1em}\selectfont\mfallback{∩}}}

% ---- 括号/开区间（图形层，不是字形）----
% 一律 \draw ... arc[start angle=, end angle=, radius=]：
%   · 端点记号背离它所界的区间；
%   · 左端/右端分别 10× 特写核对；
%   · 闭区间压在轴上时竖脊右移半个线宽；
%   · 同一张图同类记号同一种写法。

% 依赖：fontspec（已 \setsansfont{TeX Gyre Heros}）、graphicx 不需要
% 规则见同目录《符号验收清单.md》
% =====================================================================

% ---- 字体角色（全部走 sans = TeX Gyre Heros）----
\newcommand{\SIt}[1]{\textsf{\textit{#1}}}          % 斜体变量
\newcommand{\SBf}[1]{\textsf{\textbf{#1}}}          % 粗体
\newcommand{\SBfIt}[1]{\textsf{\textbf{\textit{#1}}}} % 粗斜体变量
\newcommand{\SNum}[1]{\textsf{\textup{#1}}}         % 直立数字

% ---- 上下标：字号/位置写死，避免数学样式与手调混用 ----
% 审计结论（2026-10-20）：源图下标/底字高比在 0.77–0.92em，0.72em 偏小。
% 默认 0.80em；仍需按源图逐图标定（44.4/44.5 约 0.85–0.92，39/41/46 约 0.77–0.80）。
\newcommand{\FSup}[2]{#1\raisebox{0.30em}{\fontsize{0.72em}{0.82em}\selectfont #2}}
\newcommand{\FSub}[2]{#1\raisebox{-0.30em}{\fontsize{0.80em}{0.92em}\selectfont #2}}

% ---- 分数：数字直立、方头粗杆；杆宽 = max(分子宽,分母宽) + 2pt ----
% 杆厚 4pt（配合 20–26pt 标签）；小字号图可局部调 0.17em
% 2026-10-20 修正：原 \shortstack 负 kern 不随字号缩放（23pt 下总高仅 30pt、杆贴字），
% 改为 offinterlineskip + em 定间距：上间隙 0.12em、下间隙 0.13em
% （37.2 源图 23pt 实测四标签上/下间隙均值 2.6/4.0pt；取近平衡且总高保持 42–43pt）。
\newcommand{\FFrac}[2]{%
  \begingroup
  \setbox0=\hbox{#1}\setbox2=\hbox{#2}%
  \dimen0=\wd0 \ifdim\wd2>\dimen0 \dimen0=\wd2 \fi
  \advance\dimen0 by 2pt
  \leavevmode\vbox{\offinterlineskip
    \hbox to \dimen0{\hss\box0\hss}%
    \kern0.12em
    \hbox{\rule{\dimen0}{4pt}}%
    \kern0.13em
    \hbox to \dimen0{\hss\box2\hss}%
  }%
  \endgroup}

% ---- 真减号（审计：源图减号 15–18px 宽、3–5px 厚；Heros 常规字重偏细）----
% 可选参数 = 宽度（em）：正文 0.58；上标内用 \FMinus[0.72]
\newcommand{\FMinus}[1][0.58]{\rule[0.30em]{#1em}{0.13em}}

% ---- 直撇（审计：44.2 用弯引号、44.3 用直撇，必须统一为直楔撇）----
\newcommand{\FPrime}{\textsf{\textquotesingle}}

% ---- 乘号（审计：46.2 的 × 比源小 25–30%，标签内需按源校准）----
% 用 \FTimes 时在 node 内与底字同字号，并在目检时量宽高。


% ---- 斜杠分数（审计 A：35.1 的 r/3 用 38pt 斜杠手搓）----
\newcommand{\FSlashFrac}[2]{#1\raisebox{-0.18em}{\fontsize{1.35em}{1.35em}\selectfont\textsf{/}}\kern-0.06em #2}

% ---- 交集号 ∩（Heros 无此字，固定回退字号/基线，避免逐图不同）----
\newcommand{\FIntersect}{\raisebox{0.02em}{\fontsize{0.95em}{1em}\selectfont\mfallback{∩}}}

% ---- 括号/开区间（图形层，不是字形）----
% 一律 \draw ... arc[start angle=, end angle=, radius=]：
%   · 端点记号背离它所界的区间；
%   · 左端/右端分别 10× 特写核对；
%   · 闭区间压在轴上时竖脊右移半个线宽；
%   · 同一张图同类记号同一种写法。

\begin{document}
\begin{tikzpicture}[x=1pt, y=-1pt, line cap=round, line join=round]

  % ---- labels first (one node per label, white background) ----
  \node[anchor=center,inner sep=0pt,fill=white,font=\fontsize{30}{36}\selectfont] at (516.5,33)     {\SBfIt{R}};
  \node[anchor=center,inner sep=0pt,fill=white,font=\fontsize{30}{36}\selectfont] at (531,139.5)     {\SBfIt{r}};
  \node[anchor=center,inner sep=0pt,fill=white,font=\fontsize{30}{36}\selectfont] at (537.5,263.5)   {\FSub{\SBfIt{l}}{\SBf{3}}};
  \node[anchor=center,inner sep=0pt,fill=white,font=\fontsize{30}{36}\selectfont] at (543,379.5)     {\FSlashFrac{\SBfIt{r}}{\hspace{1.5pt}\SBf{3}}};
  \node[anchor=center,inner sep=0pt,fill=white,font=\fontsize{30}{36}\selectfont] at (536,482.5)     {\FSub{\SBfIt{l}}{\SBf{2}}};
  \node[anchor=center,inner sep=0pt,fill=white,font=\fontsize{30}{36}\selectfont] at (470.5,717.5)   {\FMinus\hspace{10pt}\FSlashFrac{\SBfIt{r}}{\hspace{1.5pt}\SBf{3}}};
  \node[anchor=center,inner sep=0pt,fill=white,font=\fontsize{30}{36}\selectfont] at (533,835.5)     {\FSub{\SBfIt{l}}{\SBf{1}}};
  \node[anchor=center,inner sep=0pt,fill=white,font=\fontsize{30}{36}\selectfont] at (480.5,1002.5)  {\FMinus\hspace{10pt}\SBfIt{r}};
  \node[anchor=center,inner sep=0pt,fill=white,font=\fontsize{30}{36}\selectfont] at (273,213)       {\SBfIt{f}};
  \node[anchor=center,inner sep=0pt,fill=white,font=\fontsize{30}{36}\selectfont] at (104,470)       {\SBfIt{g}};
  \node[anchor=center,inner sep=0pt,fill=white,font=\fontsize{30}{36}\selectfont] at (278,605.5)     {\SBfIt{C}};
  \node[anchor=center,inner sep=0pt,fill=white,font=\fontsize{30}{36}\selectfont] at (612,604.5)     {\SBfIt{B}};
  \node[anchor=center,inner sep=0pt,fill=white,font=\fontsize{30}{36}\selectfont] at (1049,550)      {\SBfIt{X}};

  % ---- four dashed level lines ----
  \draw[line width=4.4, dash pattern=on 24.6pt off 15.4pt] (26,117.5) -- (1028,117.5);
  \draw[line width=4.4, dash pattern=on 24.6pt off 15.4pt] (26,401.5) -- (1028,401.5);
  \draw[line width=4.4, dash pattern=on 24.6pt off 15.4pt] (26,690.5) -- (1028,690.5);
  \draw[line width=4.4, dash pattern=on 24.6pt off 15.4pt] (26,973.5) -- (1028,973.5);

  % ---- axes ----
  \draw[line width=5]   (24,547.5) -- (1026,547.5);
  \draw[line width=4.5] (513,54)   -- (513,1032);

  % ---- six ticks on x-axis ----
  \draw[line width=4] (143.5,527) -- (143.5,563);
  \draw[line width=4] (194.5,527) -- (194.5,563);
  \draw[line width=4] (361.5,527) -- (361.5,563);
  \draw[line width=4] (533.5,527) -- (533.5,563);
  \draw[line width=4] (691.5,527) -- (691.5,563);
  \draw[line width=4] (875,527)   -- (875,563);

  % ---- braces C and B (one continuous path each) ----
  \draw[line width=5.5] (199,566) .. controls (199,572) and (202,574.5) .. (208,574.5)
    -- (264,574.5) .. controls (271,574.5) and (274,577) .. (275.5,585)
    .. controls (277,577) and (280,574.5) .. (287,574.5)
    -- (348,574.5) .. controls (354,574.5) and (357,572) .. (357,566);
  \draw[line width=5.5] (536,566) .. controls (536,572) and (539,574.5) .. (545,574.5)
    -- (599,574.5) .. controls (606,574.5) and (609,577) .. (610,585)
    .. controls (611,577) and (614,574.5) .. (621,574.5)
    -- (683,574.5) .. controls (689,574.5) and (692,572) .. (692,566);

  % ---- dashed g：左段自由曲线 → 贴合 r/3 的“直线虚线”段 → 右段自由曲线 ----
  \draw[line width=6, dash pattern=on 7pt off 12pt] plot[smooth] coordinates {
    (24,457) (53,454) (75,451) (97,446) (117,441) (138,434) (158,426)
    (180,414) (195,407) (205,401.5)
  };
  % 第一段“贴合”：直线虚线效果（g 的虚线节奏，画在 y=401.5 上），不是自由曲线
  \draw[line width=6, dash pattern=on 7pt off 12pt] (205,401.5) -- (345,401.5);
  % 中段：从 r/3 下到 -r/3
  \draw[line width=6, dash pattern=on 7pt off 12pt] plot[smooth] coordinates {
    (345,401.5) (360,408) (370,415) (380,424) (390,442) (400,456) (410,468)
    (420,476) (430,484) (440,493) (450,500) (455,499) (465,505) (471,512)
    (478,520) (486,529) (492,541) (498,554) (503,568) (512,590) (520,611)
    (526,630) (531,652) (535,670) (538,687) (540,690.5)
  };
  % 第二段“贴合”：沿 -r/3 的直线虚线（不再向下兜一个自由曲线）
  \draw[line width=6, dash pattern=on 7pt off 12pt] (540,690.5) -- (700,690.5);
  % 右段：从 -r/3 起身的可见曲线
  \draw[line width=6, dash pattern=on 7pt off 12pt] plot[smooth] coordinates {
    (700,690.5) (705,676) (716,660) (729,643) (743,626) (755,618) (762,615)
    (770,617) (779,623) (786,631) (793,639) (805,641) (814,641) (822,636)
    (832,629) (848,614) (865,599) (882,587) (901,579) (912,577) (922,578)
    (932,581) (942,586) (952,592) (960,600) (968,608) (974,617) (982,626)
    (989,634) (998,641) (1006,646) (1015,650) (1024,651)};

  % ---- solid function f (one continuous command, on top of g) ----
  \draw[line width=6] plot[smooth] coordinates {
    (146,616) (150,598) (155,566) (160,535) (165,505) (170,481) (175,461)
    (180,441) (185,421) (191,403) (196,377) (205,344) (215,312) (225,284)
    (235,261) (245,246) (255,235) (265,226) (275,222) (285,218) (293,217)
    (301,219) (309,222) (317,228) (325,237) (333,248) (341,266) (349,296)
    (355,346) (359,377) (362,403) (366,425) (370,452) (375,466) (380,474)
    (384,478) (389,474) (395,464) (400,453) (405,444) (410,437) (414,434)
    (418,437) (422,447) (426,458) (430,477) (435,506) (438,527) (441,548)
    (445,570) (450,585) (455,606) (465,618) (475,626) (485,632) (495,636)
    (505,646) (512,657) (520,666) (528,678) (536,690) (540,705) (545,716)
    (550,729) (555,744) (560,762) (565,788) (570,810) (575,820) (580,821)
    (585,818) (590,810) (595,798) (600,786) (605,782) (610,785) (615,806)
    (620,844) (625,877) (630,896) (635,901) (640,902) (645,898) (650,890)
    (655,880) (660,864) (665,842) (675,794) (685,734) (695,669) (700,641)
    (705,612) (710,582) (714,556) (718,534) (722,516) (726,500) (730,488)
    (735,473) (740,462) (745,456) (750,448) (755,446) (762,444) (770,445)
    (780,449) (790,455) (800,465) (810,478) (820,490) (830,497) (838,498)
    (845,494) (852,489) (860,478) (866,466) (871,452) (875,442) (877,437)};

  % ---- endpoint dots of f ----
  \fill (146,616) circle (9);
  \fill (877,437) circle (9);

  % ---- diamonds where R-axis meets each level line ----
  \fill (505,117.5) -- (514.5,109.5) -- (524,117.5) -- (514.5,125.5) -- cycle;
  \fill (505,401.5) -- (514.5,393.5) -- (524,401.5) -- (514.5,409.5) -- cycle;
  \fill (505,690.5) -- (514.5,682.5) -- (524,690.5) -- (514.5,698.5) -- cycle;
  \fill (505,973.5) -- (514.5,965.5) -- (524,973.5) -- (514.5,981.5) -- cycle;

  \pgfresetboundingbox
  \path[use as bounding box] (0,0) rectangle (1072,1045);
\end{tikzpicture}
\end{document}
